\section{Are responses to conspecific, heterospecific, and predator-associated sounds innate or learned?}
\label{sec:q2}

% Part I — Recognition and ecological meaning (Q2).
% See docs/OUTLINE.md — Q2.
% Draft 1. Style model: Magrath (agents/STYLE.md).

% VERIFY: MagrathEtAl2015 (not full-text checked): unlearned generalization and learning as non-exclusive routes to recognition.
A sound that predicts danger in nature need not be recognized without experience. Responses to ecologically relevant sounds could rest on unlearned recognition of species-typical structure, on acoustic properties that capture attention regardless of source, or on learning, and these alternatives are not mutually exclusive \citep{MagrathEtAl2015}. A response by a naïve receiver is consistent with innate recognition only if acoustic salience has been excluded, and a response by an experienced receiver does not show that learning was required. The rodent evidence fits neither side of a simple dichotomy. We adopt as a working model a sequence of intrinsic acoustic bias, developmental and associative calibration, and ecological experience, and evaluate it against two cases in which the evidence is most direct: avian predator calls and rat 22-kHz calls.

\subsection{Predator calls}
\label{sec:q2-predator}

Laboratory tests of predator-call recognition in predator-naïve rodents disagree. In male DBA/2 mice, calls of barn and tawny owls played at 60--70~dB after a baseline of human voice increased self-oriented, call-oriented, and defensive behavior, and the tawny owl call roughly halved food intake \citep{HendrieNeill1991}. The response was not confined to predators: kestrel and gull calls also increased defense, and the authors attributed the gull effect to attention. Using C57BL/6J mice, Sprague-Dawley rats, and pallid gerbils (\textit{Gerbillus perpallidus}), \citet{KindermannEtAl2009} found no defensive response to owl or raptor calls relative to non-predator calls, with enough power to detect a 12\% change in freezing. In rats, all calls supported fear conditioning after ten pairings without an advantage for predator calls. The two studies differ in species and strain, sound level, baseline procedure, and endpoint, and Kindermann and colleagues report that the Hendrie group later found no fear in untreated DBA/2 mice. Neither result therefore settles whether laboratory rodents carry an unlearned representation of predator calls; the positive result is the strongest challenge to a general negative claim, but it is fragile and may be specific to strain and paradigm.

Rodents with field experience respond more consistently, but their responses differ in form. Wild-caught voles, dormice, and jerboas changed their behavior to owl calls, whereas spiny mice (\textit{Acomys}) responded less \citep{HendrieEtAl1998}. In a related comparison, spiny mice showed an endocrine response to owl calls without the freezing or flight that voles expressed \citep{EilamEtAl1999} (\secref{sec:q4}).
% VERIFY: HendrieEtAl1998 and EilamEtAl1999 (not full-text checked): species differences in behavioral and endocrine responses to owl calls.
% VERIFY: HernandezEtAl2021 (not full-text checked): predator- and species-specific changes; moonlight and cover effects.
Free-living rodents of three genera (\textit{Abrothrix}, \textit{Oligoryzomys}, and \textit{Rattus}) altered feeding, movement, and vigilance in response to owl calls in ways that depended on both the owl and the rodent species, and moonlight and vegetation cover modified these effects \citep{HernandezEtAl2021}. These studies show that predator calls are meaningful to experienced receivers, but because wild-caught animals differ from laboratory strains in both experience and genetic background, they cannot separate learning from population differences in predisposition.

\subsection{Rat 22-kHz calls}
\label{sec:q2-22khz}

Conspecific 22-kHz calls in rats provide the most detailed laboratory case, and here experience dependence is well established for one endpoint. Naïve rats did not show call-specific defensive behavior to natural 22-kHz calls, although stimuli in the 22-kHz frequency range, pooled, increased freezing \citep{EndresEtAl2007}, and naïve rats did not freeze to 22-kHz calls while their amygdala neurons distinguished call types, largely by frequency \citep{ParsanaEtAl2012a} (\secref{sec:q3}). When the 22-kHz calls of a distressed partner were played to receivers, only receivers with prior aversive experience froze, and the response required the auditory thalamus \citep{KimEtAl2010}. Because the partner also displayed other fear behaviors, calls were not fully isolated from other cues in that design. Endres and colleagues placed the frequency-range advantage in the persistence of learned fear rather than in faster acquisition, which suggests that the 22-kHz band carries some intrinsic salience that facilitates retention of an acquired meaning. Selective freezing is not the only measure, however: natural 22-kHz calls slowed locomotion and heart rate without prior fear conditioning \citep{OlszynskiEtAl2020}. The claim that naïve rats do not respond to 22-kHz calls is therefore too strong; what depends on experience is selective freezing.

The mechanism of that dependence is unresolved. \citet{ParsanaEtAl2012b} proposed that rats learn the meaning of 22-kHz calls by hearing their own calls during aversive experience.
% VERIFY: ParsanaEtAl2012b (not full-text checked): autoconditioning proposal.
\citet{CalubEtAl2018} tested this by transecting the recurrent laryngeal nerve, which prevented rats from vocalizing during foot shock; call-evoked freezing was still acquired. Experience dependence survives the test, but the strict self-vocalization account does not, and what the receiver associates with the calls remains unknown.

\subsection{Development and calibration}
\label{sec:q2-development}

Developmental studies suggest that physiological discrimination can precede behavioral expression. Juvenile Belding's ground squirrels (\textit{Urocitellus beldingi}) discriminated alarm calls by heart rate before they discriminated them behaviorally \citep{Mateo1996}, and in this species trills associated with terrestrial predators raised cortisol whereas whistles associated with aerial predators were linked to slowed heart rate \citep{Mateo2010}.
% VERIFY: Mateo1996 and Mateo2010 (not full-text checked): ontogeny of heart-rate discrimination; call-type-specific physiological responses.
In laboratory mice, maternal experience calibrates cortical and behavioral tuning to pup calls on top of an intrinsic bias toward species-typical temporal structure \citep{SchiavoEtAl2020}, and association cortex contributes to this experience-dependent change \citep{TasakaEtAl2020}.
% VERIFY: TasakaEtAl2020 (not full-text checked): temporal association cortex role in auditory-driven maternal plasticity.
These systems differ in valence and context, but both fit a sequence in which an initial bias is refined by experience, and both show that the absence of an overt response in young or naïve animals does not imply the absence of discrimination.

\subsection{Eavesdropping}
\label{sec:q2-eavesdropping}

Responses to heterospecific alarm calls offer a direct way to separate learning from acoustic generalization, because the receiver's experience with the calling species can vary while the acoustic structure is known.
% VERIFY: Shriner1998, AschemeierMaher2011, WatermanMai2020 (not full-text checked): eavesdropping results.
Yellow-bellied marmots and golden-mantled ground squirrels (\textit{Callospermophilus lateralis}) distinguished alarm calls from control birdsong and responded similarly to conspecific and heterospecific alarms \citep{Shriner1998}. Woodchucks (\textit{Marmota monax}) increased vigilance to both conspecific and heterospecific alarms, more strongly to conspecific calls \citep{AschemeierMaher2011}, and Cape ground squirrels (\textit{Xerus inauris}) were as vigilant to crowned lapwing alarms as to conspecific alarms \citep{WatermanMai2020}. Eavesdropping on heterospecific alarms is thus well documented in sciurids, but these results do not reveal how it arises.

The most diagnostic test comes from eastern gray squirrels. Squirrels responded more to the alarm calls of a familiar sympatric bird, and of an unfamiliar allopatric bird with acoustically similar calls, than to the acoustically dissimilar alarm of an unfamiliar species \citep{GetschowEtAl2013}. Both familiarity and acoustic similarity therefore contribute. The contrast is weaker than it first appears, because squirrels did not distinguish the familiar bird's song from its alarm, so the response may reflect familiarity with the species rather than with its alarm calls specifically. Receivers also use heterospecific sounds as evidence of safety: after a hawk call, gray squirrels that heard ordinary bird chatter relaxed vigilance faster than squirrels that heard ambient noise \citep{LillyEtAl2019}.
% VERIFY: LillyEtAl2019 (not full-text checked): faster relaxation of vigilance with bird chatter than with ambient noise.
% VERIFY: SchmidtEtAl2008 (not full-text checked): chipmunk foraging responses to titmouse mobbing calls versus hawk calls.
An indirect cue can also outweigh a direct predator cue, as when chipmunks changed their foraging more in response to titmouse mobbing calls than to hawk calls (\secref{sec:q1}; \citealp{SchmidtEtAl2008}).

\subsection{When should recognition be innate?}
\label{sec:q2-function}

% VERIFY: Blumstein2006 (not full-text checked): persistence of antipredator traits after loss of one predator when others maintain the same defenses.
Functional theory predicts when unlearned recognition should be favored. Extending error management \citep{Nesse2005}, we expect a receiver to respond without experience to cues associated with predators whose first attack is likely to be fatal and whose acoustic signatures are predictable across generations. Learning should be favored when the predator community, or the reliability of heterospecific informants, varies in space and time. The multipredator hypothesis adds that recognition of a particular predator may persist after that predator is lost if other predators maintain selection on the same defenses \citep{Blumstein2006}. These predictions are testable with comparisons between rodent populations that do and do not coexist with a given predator, using the same calls under the same playback conditions; we found no such acoustic test in rodents. Current evidence supports a narrower conclusion: for rat 22-kHz calls, a frequency band with some intrinsic salience acquires the capacity to evoke selective freezing through experience; for predator calls, laboratory evidence of unlearned recognition is inconsistent across strains and species; and for heterospecific alarms, learning and acoustic generalization both contribute. Where a sound's significance is intrinsic and where it is acquired depends on the system, and in most systems the relevant experiments have not separated the two.

\tabref{tab:summary-q2} consolidates these claims for later condensation or figure drafting.

% Intermediate summary matrix for Q2. May later be condensed, moved to an appendix, or redrawn as a figure.
\begin{table}[tbp]
\centering
\footnotesize
\singlespacing
\caption{Working summary for Q2: innate bias versus experience. Working model: intrinsic acoustic bias $\rightarrow$ developmental and associative calibration $\rightarrow$ ecological experience.}
\label{tab:summary-q2}
\begin{tabular}{@{}p{1.3in}p{1.45in}p{2.55in}p{0.85in}@{}}
\toprule
\textbf{System} & \textbf{Na\"ive / early} & \textbf{With experience} & \textbf{Status} \\
\midrule
Avian predator calls & Conflict: defensive in one inbred mouse strain; null in mice, rats, gerbils & Field-experienced rodents respond, form varies by species & Unresolved \\
\addlinespace[2pt]
Rat 22-kHz calls & No selective freezing; frequency-band and neural discrimination present & Selective freezing after aversive experience; self-vocalization not required & Strong \\
\addlinespace[2pt]
Development (sciurids) & Heart-rate discrimination precedes behavior & Call-type physiology maps onto predator class & Supporting \\
\addlinespace[2pt]
Pup calls (mice) & Bias to species-typical intervals & Maternal experience broadens cortical and behavioral tuning & Strong \\
\addlinespace[2pt]
Heterospecific alarms & --- & Eavesdropping widespread; familiarity and acoustic similarity both contribute; safety cues also used & Strong \\
\addlinespace[2pt]
Multipredator / relaxed selection & --- & Framing only; no rodent acoustic population contrast found & Open \\
\bottomrule
\end{tabular}
\end{table}

